\documentclass[a4paper,12pt]{article}
\usepackage[margin=1in]{geometry}
\usepackage{graphicx}
\usepackage{fancyhdr}
\usepackage{setspace}
\usepackage{titlesec}
\usepackage{longtable}
\usepackage{array}
\usepackage{caption}
\usepackage{amsmath}
\usepackage{amssymb}
\usepackage{xcolor}
\usepackage{hyperref}
\usepackage{enumitem}

% Page number bottom center
\pagestyle{fancy}
\fancyhf{}
\fancyfoot[C]{\thepage}
\renewcommand{\headrulewidth}{0pt}

% Custom colors
\definecolor{questioncolor}{RGB}{0, 51, 102}
\definecolor{answercolor}{RGB}{0, 100, 0}

% Custom commands for Q&A
\newcommand{\vivaQ}[1]{\textcolor{questioncolor}{\textbf{Q: #1}}}
\newcommand{\vivaA}[1]{\textcolor{answercolor}{\textbf{A:}} #1}

\begin{document}

% ----------------- Title Page -----------------
\begin{titlepage}
    \centering
    {\Large \textbf{National University of Modern Languages, Islamabad}} \par
    \vspace{2cm}
    {\Huge \bfseries Viva Preparation Guide \par}
    \vspace{1cm}
    {\Large \textit{Real-Time Speech Recognition System} \par}
    \vspace{0.5cm}
    {\large \textit{Questions \& Answers for Oral Examination} \par}
    \vspace{3cm}
    {\large \textbf{Student:} Junaid Asif} \par
    \vspace{0.3cm}
    {\large \textbf{Roll No:} BSAI-144} \par
    \vspace{0.3cm}
    {\large \textbf{Course:} Speech Processing} \par
    \vfill
    {\large December 2025}
\end{titlepage}

\setstretch{1.3}

\tableofcontents
\newpage

% =============================================================================
\section{Project Overview Questions}
% =============================================================================

\vivaQ{What is the main objective of your project?}

\vivaA{The main objective is to design and implement a Python-based system that can recognize spoken digits (zero to nine) in real-time from a live microphone audio stream, displaying the recognized command with minimal latency and confidence percentage.}

\vspace{0.5cm}

\vivaQ{Why did you choose Python instead of MATLAB?}

\vivaA{Python was chosen because:
\begin{itemize}
    \item Free and open-source (no license required)
    \item Rich ecosystem of libraries (TensorFlow, Librosa, NumPy)
    \item Better deployment options for real-world applications
    \item Larger community support for deep learning
    \item Cross-platform compatibility
\end{itemize}}

\vspace{0.5cm}

\vivaQ{What are the main components/modules of your system?}

\vivaA{The system consists of five main modules:
\begin{enumerate}
    \item \textbf{config.py} - Configuration settings and parameters
    \item \textbf{record\_dataset.py} - Audio recording and dataset creation
    \item \textbf{feature\_extraction.py} - MFCC feature extraction
    \item \textbf{train\_model.py} - CNN model training
    \item \textbf{realtime\_recognition.py} - Live speech recognition
\end{enumerate}}

\vspace{0.5cm}

\vivaQ{What is the workflow/pipeline of your system?}

\vivaA{The pipeline is:
\begin{enumerate}
    \item Audio captured from microphone
    \item Voice Activity Detection (VAD) checks for speech
    \item MFCC features extracted from audio
    \item Features normalized using z-score
    \item CNN model performs classification
    \item Predicted digit displayed with confidence score
\end{enumerate}}

% =============================================================================
\section{Audio Processing Questions}
% =============================================================================

\vivaQ{What sampling rate did you use and why?}

\vivaA{I used \textbf{16,000 Hz (16 kHz)} because:
\begin{itemize}
    \item Human speech contains frequencies up to ~8 kHz
    \item By Nyquist theorem: $f_s \geq 2 \times f_{max}$, so 16 kHz captures all speech frequencies
    \item It's a standard rate for speech recognition systems
    \item Balances quality and computational efficiency
\end{itemize}}

\vspace{0.5cm}

\vivaQ{What is the Nyquist theorem?}

\vivaA{The Nyquist theorem states that to accurately reconstruct a signal, the sampling frequency must be at least twice the maximum frequency in the signal:
\[f_s \geq 2 \times f_{max}\]
For speech with maximum frequency of 8 kHz, we need at least 16 kHz sampling rate.}

\vspace{0.5cm}

\vivaQ{Why did you use mono audio instead of stereo?}

\vivaA{Mono audio was used because:
\begin{itemize}
    \item Speech recognition doesn't benefit from stereo information
    \item Reduces data size by half
    \item Simplifies processing pipeline
    \item Most microphones provide mono input anyway
\end{itemize}}

\vspace{0.5cm}

\vivaQ{What is the duration of each audio sample?}

\vivaA{Each audio sample is \textbf{1 second} long, which is sufficient for isolated digit utterances. This duration:
\begin{itemize}
    \item Captures complete digit pronunciation
    \item Creates fixed-length inputs for the neural network
    \item Provides consistent feature dimensions
\end{itemize}}

% =============================================================================
\section{MFCC Questions (Most Important)}
% =============================================================================

\vivaQ{What is MFCC? Explain in detail.}

\vivaA{MFCC stands for \textbf{Mel-Frequency Cepstral Coefficients}. They are features that represent the short-term power spectrum of sound in a way that mimics human auditory perception.

MFCCs capture the shape of the vocal tract, which is responsible for creating different sounds. They are widely used in speech and audio processing because they:
\begin{itemize}
    \item Compress spectral information efficiently
    \item Are robust to noise
    \item Capture phonetically important characteristics
\end{itemize}}

\vspace{0.5cm}

\vivaQ{Explain the complete MFCC extraction process step by step.}

\vivaA{The MFCC extraction process involves 7 steps:

\textbf{Step 1: Pre-emphasis}
\begin{itemize}
    \item Boosts high frequencies to balance the spectrum
    \item Formula: $y[n] = x[n] - \alpha \cdot x[n-1]$, where $\alpha \approx 0.97$
    \item Compensates for the natural roll-off in speech spectrum
\end{itemize}

\textbf{Step 2: Framing}
\begin{itemize}
    \item Divides signal into short overlapping frames (20-40ms)
    \item Speech is quasi-stationary within short frames
    \item Overlap prevents loss of information at boundaries
\end{itemize}

\textbf{Step 3: Windowing}
\begin{itemize}
    \item Applies Hamming window to each frame
    \item Formula: $w[n] = 0.54 - 0.46 \cos\left(\frac{2\pi n}{N-1}\right)$
    \item Reduces spectral leakage in FFT
\end{itemize}

\textbf{Step 4: FFT (Fast Fourier Transform)}
\begin{itemize}
    \item Converts time domain to frequency domain
    \item Computes magnitude spectrum $|X[k]|$
\end{itemize}

\textbf{Step 5: Mel Filter Bank}
\begin{itemize}
    \item Applies triangular filters on mel scale
    \item Mel scale: $m = 2595 \cdot \log_{10}(1 + f/700)$
    \item Mimics human ear's frequency resolution
\end{itemize}

\textbf{Step 6: Log Compression}
\begin{itemize}
    \item Takes logarithm of filter bank energies
    \item Mimics human loudness perception
    \item Reduces dynamic range
\end{itemize}

\textbf{Step 7: DCT (Discrete Cosine Transform)}
\begin{itemize}
    \item Decorrelates the log filter bank energies
    \item Produces final MFCC coefficients
    \item Typically keep first 12-13 coefficients
\end{itemize}}

\vspace{0.5cm}

\vivaQ{What is the Mel scale? Why is it used?}

\vivaA{The Mel scale is a perceptual scale of pitches that mimics how humans perceive frequency differences.

\textbf{Formula:}
\[m = 2595 \cdot \log_{10}\left(1 + \frac{f}{700}\right)\]

\textbf{Why it's used:}
\begin{itemize}
    \item Human ears are more sensitive to low frequencies than high frequencies
    \item Equal intervals on mel scale sound equally spaced to humans
    \item Below 1000 Hz: approximately linear
    \item Above 1000 Hz: approximately logarithmic
\end{itemize}

\textbf{Example:}
\begin{itemize}
    \item 1000 Hz = 1000 mel
    \item 2000 Hz $\approx$ 1500 mel (not 2000 mel)
\end{itemize}}

\vspace{0.5cm}

\vivaQ{What is the difference between MFCC and Mel Spectrogram?}

\vivaA{
\begin{tabular}{|l|l|}
\hline
\textbf{Mel Spectrogram} & \textbf{MFCC} \\
\hline
Output of mel filter bank & Output after DCT on mel spectrogram \\
\hline
Correlated features & Decorrelated features \\
\hline
Higher dimensionality & Lower dimensionality \\
\hline
2D representation & 1D coefficient vector per frame \\
\hline
Used in modern deep learning & Traditional speech recognition \\
\hline
\end{tabular}}

\vspace{0.5cm}

\vivaQ{Why do you use 13 MFCC coefficients?}

\vivaA{13 coefficients are used because:
\begin{itemize}
    \item First 13 coefficients capture most vocal tract information
    \item Higher coefficients represent fine spectral details (often noise)
    \item Industry standard for speech recognition
    \item Balance between information and dimensionality
    \item The 0th coefficient represents frame energy (sometimes excluded)
\end{itemize}}

\vspace{0.5cm}

\vivaQ{What are Delta and Delta-Delta coefficients?}

\vivaA{\textbf{Delta coefficients} (also called velocity or differential coefficients) capture the rate of change of MFCCs over time - they represent \textit{how} the sound is changing.

\textbf{Delta-Delta coefficients} (also called acceleration coefficients) capture the rate of change of delta coefficients - they represent the \textit{acceleration} of spectral changes.

\textbf{Formula for Delta:}
\[\Delta c_t = \frac{\sum_{n=1}^{N} n(c_{t+n} - c_{t-n})}{2\sum_{n=1}^{N} n^2}\]

\textbf{Why they're important:}
\begin{itemize}
    \item Static MFCCs only capture instantaneous spectrum
    \item Delta captures transitional information
    \item Transitions between phonemes are important for recognition
    \item Total features: 13 + 13 + 13 = 39 coefficients
\end{itemize}}

\vspace{0.5cm}

\vivaQ{What is N\_FFT and HOP\_LENGTH?}

\vivaA{\textbf{N\_FFT (512):}
\begin{itemize}
    \item FFT window size in samples
    \item At 16 kHz: $512/16000 = 32$ ms window
    \item Determines frequency resolution
    \item Must be power of 2 for efficient FFT
\end{itemize}

\textbf{HOP\_LENGTH (160):}
\begin{itemize}
    \item Number of samples between consecutive frames
    \item At 16 kHz: $160/16000 = 10$ ms hop
    \item Controls time resolution and overlap
    \item Overlap = N\_FFT - HOP\_LENGTH = 352 samples
\end{itemize}

\textbf{Trade-off:}
\begin{itemize}
    \item Larger N\_FFT $\rightarrow$ better frequency resolution, worse time resolution
    \item Smaller HOP\_LENGTH $\rightarrow$ more frames, higher computation
\end{itemize}}

\vspace{0.5cm}

\vivaQ{What is the shape of your feature matrix and why?}

\vivaA{The feature shape is \textbf{(39, 101, 1)} where:
\begin{itemize}
    \item \textbf{39} = Number of features (13 MFCC + 13 Delta + 13 Delta-Delta)
    \item \textbf{101} = Number of time frames = $\lceil 16000/160 \rceil + 1$
    \item \textbf{1} = Single channel (like grayscale image)
\end{itemize}

This shape allows the CNN to process MFCC features as a 2D image where:
\begin{itemize}
    \item Height represents frequency/cepstral information
    \item Width represents time progression
\end{itemize}}

% =============================================================================
\section{Voice Activity Detection (VAD) Questions}
% =============================================================================

\vivaQ{What is Voice Activity Detection (VAD)?}

\vivaA{VAD is a technique to determine whether an audio segment contains speech or just silence/noise. It acts as a gate that:
\begin{itemize}
    \item Prevents processing during silence
    \item Reduces false positives
    \item Saves computational resources
    \item Improves system responsiveness
\end{itemize}}

\vspace{0.5cm}

\vivaQ{What is RMS and how is it used for VAD?}

\vivaA{RMS stands for \textbf{Root Mean Square}. It measures the average energy/loudness of an audio signal.

\textbf{Formula:}
\[RMS = \sqrt{\frac{1}{n}\sum_{i=1}^{n} x_i^2}\]

\textbf{For VAD:}
\begin{itemize}
    \item If $RMS > threshold$ (0.01) $\rightarrow$ Speech detected
    \item If $RMS < threshold$ $\rightarrow$ Silence, skip processing
\end{itemize}

\textbf{Why RMS instead of simple average:}
\begin{itemize}
    \item Audio signals oscillate around zero
    \item Simple average would be near zero
    \item Squaring makes all values positive
    \item RMS gives meaningful energy measure
\end{itemize}}

\vspace{0.5cm}

\vivaQ{What are other VAD methods?}

\vivaA{Other VAD methods include:
\begin{enumerate}
    \item \textbf{Zero Crossing Rate (ZCR)} - Count sign changes in signal
    \item \textbf{Spectral Entropy} - Measure spectral randomness
    \item \textbf{Energy in specific frequency bands} - Focus on speech frequencies
    \item \textbf{Machine Learning based} - Train classifier for speech/non-speech
    \item \textbf{WebRTC VAD} - Google's robust VAD implementation
\end{enumerate}}

% =============================================================================
\section{CNN Model Questions}
% =============================================================================

\vivaQ{Why did you choose CNN for this task?}

\vivaA{CNN (Convolutional Neural Network) was chosen because:
\begin{itemize}
    \item MFCC features can be treated as 2D images
    \item CNNs excel at learning spatial hierarchies
    \item They capture local patterns (like phonemes)
    \item Parameter sharing reduces model size
    \item Translation invariance helps with timing variations
    \item Proven effectiveness in audio classification
\end{itemize}}

\vspace{0.5cm}

\vivaQ{Explain your CNN architecture.}

\vivaA{My CNN has the following architecture:

\textbf{4 Convolutional Blocks:}
\begin{enumerate}
    \item Conv2D(32 filters) $\rightarrow$ BatchNorm $\rightarrow$ ReLU $\rightarrow$ MaxPool $\rightarrow$ Dropout(0.25)
    \item Conv2D(64 filters) $\rightarrow$ BatchNorm $\rightarrow$ ReLU $\rightarrow$ MaxPool $\rightarrow$ Dropout(0.25)
    \item Conv2D(128 filters) $\rightarrow$ BatchNorm $\rightarrow$ ReLU $\rightarrow$ MaxPool $\rightarrow$ Dropout(0.25)
    \item Conv2D(256 filters) $\rightarrow$ BatchNorm $\rightarrow$ ReLU
\end{enumerate}

\textbf{Classification Head:}
\begin{enumerate}
    \item Global Average Pooling
    \item Dense(128) $\rightarrow$ BatchNorm $\rightarrow$ ReLU $\rightarrow$ Dropout(0.5)
    \item Dense(64) $\rightarrow$ BatchNorm $\rightarrow$ ReLU $\rightarrow$ Dropout(0.5)
    \item Dense(10) $\rightarrow$ Softmax
\end{enumerate}

Total parameters: ~430,000}

\vspace{0.5cm}

\vivaQ{What is the purpose of each layer type?}

\vivaA{
\textbf{Conv2D:}
\begin{itemize}
    \item Learns local patterns using learnable filters
    \item Extracts features like edges, textures, phonemes
\end{itemize}

\textbf{BatchNormalization:}
\begin{itemize}
    \item Normalizes layer outputs
    \item Speeds up training
    \item Reduces internal covariate shift
    \item Acts as regularization
\end{itemize}

\textbf{ReLU Activation:}
\begin{itemize}
    \item $f(x) = max(0, x)$
    \item Introduces non-linearity
    \item Solves vanishing gradient problem
\end{itemize}

\textbf{MaxPooling:}
\begin{itemize}
    \item Reduces spatial dimensions
    \item Provides translation invariance
    \item Reduces computation
\end{itemize}

\textbf{Dropout:}
\begin{itemize}
    \item Randomly zeros neurons during training
    \item Prevents overfitting
    \item Forces network to learn redundant representations
\end{itemize}

\textbf{GlobalAveragePooling:}
\begin{itemize}
    \item Averages each feature map to single value
    \item Reduces parameters vs Flatten
    \item More robust to spatial variations
\end{itemize}

\textbf{Dense:}
\begin{itemize}
    \item Fully connected layers
    \item Combines features for final classification
\end{itemize}

\textbf{Softmax:}
\begin{itemize}
    \item Converts logits to probabilities
    \item Outputs sum to 1
    \item $\sigma(z_i) = \frac{e^{z_i}}{\sum_j e^{z_j}}$
\end{itemize}}

\vspace{0.5cm}

\vivaQ{What loss function did you use and why?}

\vivaA{I used \textbf{Sparse Categorical Cross-Entropy} because:
\begin{itemize}
    \item Multi-class classification problem (10 classes)
    \item Labels are integers (0-9), not one-hot encoded
    \item "Sparse" version is memory efficient
\end{itemize}

\textbf{Formula:}
\[L = -\sum_{i} y_i \log(\hat{y}_i)\]

For sparse version, only the true class contributes to loss.}

\vspace{0.5cm}

\vivaQ{What optimizer did you use?}

\vivaA{I used \textbf{Adam optimizer} (Adaptive Moment Estimation) because:
\begin{itemize}
    \item Combines benefits of AdaGrad and RMSprop
    \item Adaptive learning rates for each parameter
    \item Works well with default hyperparameters
    \item Efficient with sparse gradients
    \item Learning rate: 0.001 (default)
\end{itemize}}

\vspace{0.5cm}

\vivaQ{What callbacks did you use during training?}

\vivaA{Three callbacks were used:

\textbf{1. EarlyStopping:}
\begin{itemize}
    \item Monitors validation loss
    \item Stops if no improvement for 10 epochs
    \item Restores best weights
    \item Prevents overfitting
\end{itemize}

\textbf{2. ReduceLROnPlateau:}
\begin{itemize}
    \item Reduces learning rate when stuck
    \item Factor: 0.5 (halves LR)
    \item Patience: 5 epochs
    \item Helps escape local minima
\end{itemize}

\textbf{3. ModelCheckpoint:}
\begin{itemize}
    \item Saves best model based on validation accuracy
    \item Ensures best model is preserved
\end{itemize}}

\vspace{0.5cm}

\vivaQ{What is overfitting and how did you prevent it?}

\vivaA{\textbf{Overfitting} occurs when a model learns the training data too well, including noise, and fails to generalize to new data.

\textbf{Prevention techniques used:}
\begin{enumerate}
    \item \textbf{Dropout} (0.25 in conv layers, 0.5 in dense layers)
    \item \textbf{Batch Normalization} - regularization effect
    \item \textbf{Early Stopping} - stop before overfitting
    \item \textbf{Data Augmentation} (recommended for improvement)
    \item \textbf{Small model} relative to data size
\end{enumerate}}

% =============================================================================
\section{Training Questions}
% =============================================================================

\vivaQ{How did you split your dataset?}

\vivaA{I used an 80-20 split:
\begin{itemize}
    \item \textbf{Training set:} 80\% of data
    \item \textbf{Test/Validation set:} 20\% of data
\end{itemize}

For small datasets (< 50 samples), stratified splitting was disabled to avoid errors. With 30 samples, I got 24 training and 6 test samples.}

\vspace{0.5cm}

\vivaQ{What is stratified splitting?}

\vivaA{Stratified splitting ensures that each split has the same proportion of classes as the original dataset.

\textbf{Example:} If dataset has 10\% of class A, both train and test sets will have ~10\% of class A.

\textbf{Problem with small datasets:} With only 3 samples per class and 20\% test split, we can't ensure each class appears in test set.}

\vspace{0.5cm}

\vivaQ{What is normalization and why is it important?}

\vivaA{Normalization scales features to a standard range.

\textbf{Z-score normalization used:}
\[X_{normalized} = \frac{X - \mu}{\sigma}\]

\textbf{Why it's important:}
\begin{itemize}
    \item Neural networks work better with small input values
    \item Prevents features with large values from dominating
    \item Speeds up training convergence
    \item Required for consistent inference (same params used)
\end{itemize}}

\vspace{0.5cm}

\vivaQ{Why do you save normalization parameters?}

\vivaA{Normalization parameters (mean and std) computed on training data must be saved because:
\begin{itemize}
    \item Same parameters must be used during inference
    \item New audio must be normalized the same way
    \item Different normalization would give wrong predictions
    \item Ensures consistency between training and testing
\end{itemize}}

% =============================================================================
\section{Real-Time Recognition Questions}
% =============================================================================

\vivaQ{How does real-time recognition work?}

\vivaA{Real-time recognition uses a sliding buffer approach:

\begin{enumerate}
    \item \textbf{Audio Stream:} Continuous capture from microphone
    \item \textbf{Circular Buffer:} Maintains last 1 second of audio
    \item \textbf{Callback Function:} Called every 100ms with new audio
    \item \textbf{Buffer Update:} Old samples shifted out, new samples added
    \item \textbf{Processing Loop:} Every 100ms:
    \begin{itemize}
        \item Check VAD
        \item Extract features
        \item Run model inference
        \item Display result if confident
    \end{itemize}
\end{enumerate}}

\vspace{0.5cm}

\vivaQ{What is the confidence threshold and why is it used?}

\vivaA{The confidence threshold (0.6 or 60\%) is the minimum probability required to accept a prediction.

\textbf{Purpose:}
\begin{itemize}
    \item Reduces false positives
    \item Filters uncertain predictions
    \item Only displays confident results
\end{itemize}

\textbf{Trade-off:}
\begin{itemize}
    \item Higher threshold $\rightarrow$ fewer false positives, might miss valid speech
    \item Lower threshold $\rightarrow$ more detections, more errors
\end{itemize}}

\vspace{0.5cm}

\vivaQ{What is the cooldown period?}

\vivaA{The cooldown period (0.5 seconds) prevents repeated recognitions of the same utterance.

\textbf{Problem without cooldown:}
\begin{itemize}
    \item Same word might be in buffer for multiple cycles
    \item Would output "FIVE FIVE FIVE FIVE..."
\end{itemize}

\textbf{Solution:}
\begin{itemize}
    \item After recognition, wait 0.5 seconds before next output
    \item Allows buffer to clear previous audio
\end{itemize}}

\vspace{0.5cm}

\vivaQ{What is the latency of your system?}

\vivaA{The system latency is approximately \textbf{100-150ms}, which includes:
\begin{itemize}
    \item Audio buffer filling: ~100ms (callback block size)
    \item Feature extraction: ~10-20ms
    \item Model inference: ~20-30ms
    \item Display: negligible
\end{itemize}

This is well below human-perceptible delay (~200ms).}

% =============================================================================
\section{Technical Deep-Dive Questions}
% =============================================================================

\vivaQ{What is the Fourier Transform?}

\vivaA{The Fourier Transform decomposes a signal into its constituent frequencies.

\textbf{Continuous Fourier Transform:}
\[X(f) = \int_{-\infty}^{\infty} x(t) e^{-j2\pi ft} dt\]

\textbf{Discrete Fourier Transform (DFT):}
\[X[k] = \sum_{n=0}^{N-1} x[n] e^{-j2\pi kn/N}\]

\textbf{FFT (Fast Fourier Transform):}
\begin{itemize}
    \item Efficient algorithm for computing DFT
    \item Complexity: $O(N \log N)$ instead of $O(N^2)$
    \item Requires N to be power of 2
\end{itemize}}

\vspace{0.5cm}

\vivaQ{What is the Hamming window and why is it used?}

\vivaA{The Hamming window is applied to each frame before FFT.

\textbf{Formula:}
\[w[n] = 0.54 - 0.46 \cos\left(\frac{2\pi n}{N-1}\right)\]

\textbf{Purpose:}
\begin{itemize}
    \item Reduces spectral leakage
    \item Smooth transition at frame edges
    \item Without windowing, abrupt edges cause artifacts
\end{itemize}}

\vspace{0.5cm}

\vivaQ{What is the DCT and why is it used in MFCC?}

\vivaA{DCT (Discrete Cosine Transform) converts the log mel spectrum to cepstral coefficients.

\textbf{Purpose:}
\begin{itemize}
    \item Decorrelates the features
    \item Compresses information into fewer coefficients
    \item First few coefficients contain most energy
    \item Similar to PCA effect
\end{itemize}

\textbf{Why DCT over FFT for this step:}
\begin{itemize}
    \item Input is real-valued (log energies)
    \item DCT produces real output
    \item Better energy compaction
\end{itemize}}

\vspace{0.5cm}

\vivaQ{What is a Convolutional layer doing mathematically?}

\vivaA{A convolutional layer performs the convolution operation:
\[(f * g)[n] = \sum_{m} f[m] \cdot g[n-m]\]

In 2D for images:
\[y[i,j] = \sum_m \sum_n x[i+m, j+n] \cdot w[m,n] + b\]

\textbf{Key concepts:}
\begin{itemize}
    \item \textbf{Kernel/Filter:} Small learnable weight matrix
    \item \textbf{Stride:} Step size for sliding the kernel
    \item \textbf{Padding:} Adding zeros to preserve dimensions
    \item \textbf{Feature map:} Output of applying one filter
\end{itemize}}

% =============================================================================
\section{Troubleshooting Questions}
% =============================================================================

\vivaQ{What was the stratified split error and how did you solve it?}

\vivaA{\textbf{Error:} "test\_size = 6 should be greater or equal to number of classes = 10"

\textbf{Cause:} With 30 samples and 20\% test split = 6 test samples. Stratified split needs at least 1 sample per class (10) in test set.

\textbf{Solution:} Added conditional check:
\begin{itemize}
    \item If test samples < number of classes
    \item Use simple random split instead of stratified
\end{itemize}}

\vspace{0.5cm}

\vivaQ{How would you improve accuracy with more data?}

\vivaA{With more data, I would:
\begin{enumerate}
    \item Record 20+ samples per digit
    \item Use data augmentation (noise, pitch shift, time stretch)
    \item Enable stratified splitting
    \item Increase model complexity
    \item Add cross-validation
    \item Record from different environments
    \item Include multiple speakers
\end{enumerate}}

% =============================================================================
\section{Comparison and Alternative Questions}
% =============================================================================

\vivaQ{What are alternatives to MFCC?}

\vivaA{Alternative features include:
\begin{enumerate}
    \item \textbf{Mel Spectrogram} - Raw mel-scaled spectrogram
    \item \textbf{Filter Banks} - Before DCT step
    \item \textbf{PLP} - Perceptual Linear Prediction
    \item \textbf{Spectrograms} - Time-frequency representation
    \item \textbf{Raw Waveform} - End-to-end learning
    \item \textbf{Wav2Vec} - Self-supervised representations
\end{enumerate}}

\vspace{0.5cm}

\vivaQ{What are alternatives to CNN?}

\vivaA{Alternative models include:
\begin{enumerate}
    \item \textbf{RNN/LSTM} - Sequential processing
    \item \textbf{GRU} - Simplified LSTM
    \item \textbf{Transformer} - Attention-based
    \item \textbf{CRNN} - CNN + RNN hybrid
    \item \textbf{ResNet} - Residual connections
    \item \textbf{HMM-GMM} - Traditional approach
\end{enumerate}}

\vspace{0.5cm}

\vivaQ{How is your system different from commercial systems like Siri/Alexa?}

\vivaA{
\begin{tabular}{|l|l|l|}
\hline
\textbf{Aspect} & \textbf{My System} & \textbf{Commercial} \\
\hline
Vocabulary & 10 digits & Unlimited \\
\hline
Model size & ~430K params & Billions \\
\hline
Training data & 30 samples & Millions of hours \\
\hline
Processing & Local & Cloud-based \\
\hline
Language model & None & Advanced LM \\
\hline
Noise handling & Basic VAD & Advanced DNNs \\
\hline
\end{tabular}}

% =============================================================================
\section{Quick Reference Formulas}
% =============================================================================

\textbf{Key Formulas to Remember:}

\begin{enumerate}
    \item \textbf{Mel Scale:} $m = 2595 \cdot \log_{10}(1 + f/700)$
    
    \item \textbf{RMS:} $RMS = \sqrt{\frac{1}{n}\sum_{i=1}^{n} x_i^2}$
    
    \item \textbf{Delta:} $\Delta c_t = \frac{\sum_{n=1}^{N} n(c_{t+n} - c_{t-n})}{2\sum_{n=1}^{N} n^2}$
    
    \item \textbf{Normalization:} $X_{norm} = \frac{X - \mu}{\sigma}$
    
    \item \textbf{Softmax:} $\sigma(z_i) = \frac{e^{z_i}}{\sum_j e^{z_j}}$
    
    \item \textbf{ReLU:} $f(x) = \max(0, x)$
    
    \item \textbf{Nyquist:} $f_s \geq 2 \times f_{max}$
    
    \item \textbf{Frame duration:} $T = \frac{N_{FFT}}{f_s}$ (512/16000 = 32ms)
    
    \item \textbf{Hop duration:} $T_{hop} = \frac{HOP}{f_s}$ (160/16000 = 10ms)
\end{enumerate}

% =============================================================================
\section{Key Numbers to Remember}
% =============================================================================

\begin{tabular}{|l|l|}
\hline
\textbf{Parameter} & \textbf{Value} \\
\hline
Sample Rate & 16,000 Hz \\
\hline
Duration & 1 second \\
\hline
N\_FFT & 512 samples (32ms) \\
\hline
HOP\_LENGTH & 160 samples (10ms) \\
\hline
N\_MFCC & 13 coefficients \\
\hline
Total Features & 39 (13 × 3) \\
\hline
Time Frames & 101 \\
\hline
Input Shape & (39, 101, 1) \\
\hline
Classes & 10 (digits 0-9) \\
\hline
Model Parameters & ~430,000 \\
\hline
Confidence Threshold & 60\% \\
\hline
VAD Threshold & 0.01 RMS \\
\hline
\end{tabular}

\vspace{1cm}
\begin{center}
\textbf{Good Luck with Your Viva!}
\end{center}

\end{document}
